\documentclass{beamer}
\usepackage[utf8]{inputenc}
\usepackage[brazilian]{babel}
\usepackage{booktabs}
\usepackage{xcolor}

% Tema CERISE — mesma identidade visual dos outros slides do projeto.
\definecolor{cerisemagenta}{HTML}{A6186B}
\definecolor{ceriserosa}{HTML}{E6186D}
\definecolor{cerisecoral}{HTML}{F5642D}
\definecolor{ceriseroxo}{HTML}{5B1A6B}

\setbeamercolor{title}{fg=cerisemagenta}
\setbeamercolor{frametitle}{fg=white,bg=cerisemagenta}
\setbeamercolor{structure}{fg=cerisemagenta}
\setbeamercolor{block title}{fg=white,bg=ceriseroxo}
\setbeamercolor{block body}{bg=cerisemagenta!5}
\setbeamercolor{alerted text}{fg=cerisecoral}
\setbeamercolor{item}{fg=cerisemagenta}

\newenvironment{numerobloco}[1]{\begin{block}{#1}}{\end{block}}
\newenvironment{decisaobloco}[1]{\begin{alertblock}{#1}}{\end{alertblock}}

\title{Medir antes de decidir}
\subtitle{Como o Harbor ganhou instrumentação --- e o que ela revelou}
\author{Yan Santos Leite}
\date{8--13 de setembro de 2026}

\begin{document}

\frame{\titlepage}

% ─────────────────────────────────────────────────────────────────────
% ATO 0 — a tese
% ─────────────────────────────────────────────────────────────────────
\begin{frame}{A pergunta que organiza tudo}
  \begin{numerobloco}{O problema}
    Até 8 de setembro, o Harbor não conseguia responder a uma pergunta básica sobre si
    mesmo: \emph{``essa mudança melhorou o sistema?''} --- não havia trace, não havia
    custo medido, não havia comparação antes/depois.
  \end{numerobloco}

  \pause
  \textbf{O que fizemos, em duas etapas:}
  \begin{enumerate}
    \item Construímos a capacidade de medir (trace, custo por chamada, relatório de
      replicação).
    \item Usamos essa capacidade numa série de decisões reais --- e em \alert{várias
      delas} a resposta foi \emph{não}.
  \end{enumerate}

  \pause
  \begin{decisaobloco}{A tese deste documento}
    Instrumentação não serve para confirmar o que já queríamos fazer. Serve para
    descobrir o que \emph{não} devemos fazer --- e isso só tem valor se a gente aceitar
    o ``não'' quando ele aparece.
  \end{decisaobloco}
\end{frame}

% ─────────────────────────────────────────────────────────────────────
% ATO 1 — construir a régua
% ─────────────────────────────────────────────────────────────────────
\begin{frame}{Ato 1 --- Construir a régua}
  \small
  \begin{numerobloco}{Três peças novas}
    \begin{itemize}
      \item \texttt{shared/trace.py} --- esquema espelhando \texttt{gen\_ai.*}
        (OpenTelemetry) mais uma camada de proveniência (evidências, relações
        SUPPORT/CONTRADICT).
      \item \texttt{shared/ollama\_client.py} --- captura tokens e modelo em toda
        chamada. Antes: 13 pontos do código descartavam essa informação silenciosamente.
      \item \texttt{shared/rollout\_card.py} --- relatório de replicação gerado a partir
        dos traces.
    \end{itemize}
  \end{numerobloco}

  \pause
  \textbf{Por que isso vem primeiro:} sem tokens medidos, ``LangGraph é mais caro?'' é
  opinião. Sem trace, ``o agente se corrigiu?'' é impressão. Toda comparação nos atos
  seguintes depende desta etapa existir.
\end{frame}

\begin{frame}{Ato 1 --- A régua em uso: agente ReAct sobre benchmark}
  \small
  \begin{numerobloco}{Duas tarefas, com trace completo}
    \begin{center}
    \begin{tabular}{lccc}
      \toprule
      \textbf{Tarefa} & \textbf{Execuções} & \textbf{Sucesso} & \textbf{Tokens (in/out)} \\
      \midrule
      SQL (self-repair)     & 3 & 1/3 & 27.163 / 1.342 \\
      HotpotQA (replicação) & 2 & 1/2 & --- \\
      \bottomrule
    \end{tabular}
    \end{center}
    \footnotesize Amostra pequena --- valida o padrão de trace, não sustenta conclusão
    estatística sobre o modelo.
  \end{numerobloco}

  \pause
  \textbf{O que o trace revelou:} 2 das 3 execuções SQL travaram no \emph{mesmo} erro de
  sintaxe, repetido, sem o modelo usar a Observation de erro para se corrigir. Sem
  trace, isso apareceria como ``falhou'' --- com trace, sabemos \emph{por quê}.

  \pause
  \begin{decisaobloco}{Regra que ficou}
    Todo experimento com agente entrega três artefatos: código, trace e relatório de
    replicação. Um resultado sem trace não é resultado.
  \end{decisaobloco}
\end{frame}

% ─────────────────────────────────────────────────────────────────────
% ATO 2 — a régua diz sim
% ─────────────────────────────────────────────────────────────────────
\begin{frame}{Ato 2 --- Quando a régua diz sim: migração para LangGraph}
  \small
  \begin{numerobloco}{Self-repair e roteador, medidos antes e depois}
    \begin{center}
    \begin{tabular}{lcc}
      \toprule
      \textbf{Métrica} & \textbf{Antes} & \textbf{Depois} \\
      \midrule
      Roteamento (golden set, 67 perguntas) & 56/67 & 55/67$^*$ \\
      Faithfulness médio & 64\% & 64\% \\
      Hallucination Rate (nº de alucinações) & 4 & \alert{3} \\
      Divergências de rota (roteador migrado) & --- & \alert{0/67} \\
      \bottomrule
    \end{tabular}
    \end{center}
    \footnotesize $^*$ pergunta de fronteira que oscila \emph{também} no roteador
    original --- instabilidade do desempate por LLM, não regressão.
  \end{numerobloco}

  \pause
  \textbf{O que isso significou:} migração estrutural, não comportamental --- mesmos
  prompts, mesmos 11 gates, mesmo teto de 1 tentativa.

  \pause
  \begin{decisaobloco}{Decisão}
    Ambos promovidos para produção. O item de maior risco do roadmap (portar 9+ gates
    sem perder cobertura) fechou com zero divergências medidas.
  \end{decisaobloco}
\end{frame}

\begin{frame}{Ato 2 --- O que a régua não pega: bugs só visíveis em produção}
  \small
  \begin{numerobloco}{O achado}
    Os testes automatizados passavam. O dashboard real (\texttt{streamlit run app.py})
    quebrava com \texttt{ModuleNotFoundError} --- três bugs de import em cadeia.
  \end{numerobloco}

  \pause
  \textbf{Causa:} os testes sempre rodaram com o root do projeto já no \texttt{sys.path}
  por acidente. O Streamlit, rodando de dentro de \texttt{dashboard/}, não tem esse
  privilégio --- e aí o nome \texttt{nl\_to\_sql} vira ambíguo entre ``módulo solto'' e
  ``pacote'', dependendo de quem importa primeiro.

  \pause
  \begin{decisaobloco}{Lição incorporada}
    Métrica verde não substitui abrir o sistema e usar. Fix: carregar os módulos por
    caminho de arquivo explícito (\texttt{importlib}), eliminando a ambiguidade.
  \end{decisaobloco}
\end{frame}

% ─────────────────────────────────────────────────────────────────────
% ATO 3 — a régua diz não
% ─────────────────────────────────────────────────────────────────────
\begin{frame}{Ato 3 --- Quando a régua diz não (I): Agentic RAG}
  \small
  \begin{numerobloco}{A hipótese e o número}
    Um \emph{critic node} avaliaria se o texto recuperado responde à pergunta; se não,
    reformularia a query e buscaria de novo (padrão Self-RAG/CRAG).
    \begin{center}
    Recall@1 (15 perguntas): 12/15 chamada única $\to$ \alert{11/15} com critic node.
    \end{center}
  \end{numerobloco}

  \pause
  \textbf{Causa raiz investigada:} o documento certo \emph{é} recuperado --- mas o chunk
  específico às vezes não contém a resposta (está em outra seção do mesmo arquivo).
  Reformular a \emph{query} não resolve um problema de \emph{chunking}, e pode afastar do
  documento que já havia sido encontrado.

  \pause
  \begin{decisaobloco}{Decisão}
    Não promovido. Fica como protótipo medido e documentado --- com a causa raiz
    registrada, para quem retomar não repetir o mesmo caminho.
  \end{decisaobloco}
\end{frame}

\begin{frame}{Ato 3 --- Quando a régua diz não (II): Contextual Retrieval}
  \small
  \begin{numerobloco}{A hipótese e o número}
    Prependar a cada chunk um resumo gerado por LLM situando-o no documento (técnica da
    Anthropic, adaptada para Ollama local).
    \begin{center}
    \begin{tabular}{lccc}
      \toprule
      \textbf{Métrica} & \textbf{Sem} & \textbf{Com} & \textbf{Diferença} \\
      \midrule
      Recall@5    & 98,0\% & 98,0\% & idêntico \\
      Precision@5 & 50,0\% & \alert{46,4\%} & $-3,6$pp \\
      MRR         & 0,899  & 0,907  & $+0,008$ \\
      \bottomrule
    \end{tabular}
    \end{center}
  \end{numerobloco}

  \pause
  \textbf{Causa provável:} os manuais do Harbor têm 40--181 linhas. A literatura original
  mede ganho em documentos longos --- textos curtos já carregam contexto suficiente, e o
  resumo extra dilui a precisão lexical do BM25.

  \pause
  \begin{decisaobloco}{Decisão}
    Não promovido, mas mantido como flag de ablação --- vale revisitar se o corpus crescer
    com documentos longos. Diferente do Agentic RAG, aqui não há regressão de recall.
  \end{decisaobloco}
\end{frame}

\begin{frame}{Ato 3 --- Quando a régua está quebrada: calibração de top-p}
  \small
  \begin{numerobloco}{O que parecia ser}
    O benchmark multimodal mostrava top-p com desempenho ruim no corpus pequeno
    (nDCG@5 = 0,286). Hipótese inicial: limiar mal calibrado para poucos documentos.
  \end{numerobloco}

  \pause
  \textbf{O que era de verdade:} testamos os limiares 0,7, 0,85 e 0,99 --- \alert{resultado
  idêntico nos três}. O 1º estágio chama \texttt{buscar(usar\_rerank=False)}, que devolve
  \texttt{score=0.0} para todos os candidatos. Não havia massa de score para o limiar
  cortar; a variável nunca foi exercida.

  \pause
  \begin{decisaobloco}{Decisão}
    Nem sim nem não --- a medição estava medindo a coisa errada. Item de continuidade:
    dar ao 1º estágio um sinal de score real antes de tentar qualquer calibração.
  \end{decisaobloco}
\end{frame}

\begin{frame}{Ato 3 --- Consertada a régua: 0,7 vence, e uma pergunta melhor aparece (10--11/09)}
  \small
  \begin{numerobloco}{Com score real, a calibração finalmente respondeu}
    \begin{center}
    \footnotesize
    \begin{tabular}{lccc}
      \toprule
      \textbf{Limiar} & \textbf{nDCG@5} & \textbf{Recall@5} & \textbf{Custo} \\
      \midrule
      sem rerank         & 0,864          & 93\%          & 38--62s \\
      \textbf{top-p=0,7} & \textbf{0,889} & \textbf{97\%} & \textbf{$\sim$3015s} \\
      top-p=0,85         & 0,802          & 86\%          & $\sim$6734s \\
      top-p=0,99         & 0,780          & 86\%          & $\sim$7055s \\
      \bottomrule
    \end{tabular}
    \end{center}
    \footnotesize Corpus Harbor (26 imagens), score RRF real, legendas \texttt{qwen3-vl}.
  \end{numerobloco}

  \pause
  \textbf{O que isso significou:} 0,7 venceu em todas as métricas e ainda foi mais barato
  --- limiares mais permissivos incluem candidatos ruidosos e pioram o resultado, contra
  a intuição de que ``mais candidatos'' ajudaria.

  \pause
  \begin{decisaobloco}{Decisão}
    \texttt{TOP\_P\_LIMIAR} atualizado de 0,85 (chute nunca validado) para 0,7. Mas essa
    calibração usou o \emph{mesmo} VLM nos dois lados --- a pergunta que faltava responder
    é se o rerank compensa \emph{trocar} de VLM, não só ajustar o limiar.
  \end{decisaobloco}
\end{frame}

\begin{frame}{Ato 3 --- O experimento decisivo: rerank não compensa captioning ruim}
  \small
  \begin{numerobloco}{A pergunta e o teste}
    O desenho de 2 estágios (retrieval barato + rerank caro, padrão HEAVEN,
    arXiv:2510.22215) promete recuperar qualidade mesmo com um encoder fraco.
    Comparamos: (A) caption ruim (\texttt{moondream}) + rerank \texttt{top-p=0,7} vs.\
    (B) caption bom (\texttt{qwen3-vl}) sem nenhum rerank.
    \begin{center}
    \footnotesize
    \begin{tabular}{lccc}
      \toprule
      \textbf{Cenário} & \textbf{nDCG@5} & \textbf{Recall@5} & \textbf{Custo} \\
      \midrule
      A) moondream + rerank & 0,436 & 48\% & 4779s \\
      \textbf{B) qwen3-vl, sem rerank} & \textbf{0,864} & \textbf{93\%} & \textbf{11,8s} \\
      \bottomrule
    \end{tabular}
    \end{center}
  \end{numerobloco}

  \pause
  \textbf{O que isso significou:} (B) supera (A) em todas as métricas por margem larga
  (quase o dobro de nDCG/Recall) e custa \alert{$\sim$400$\times$ menos} --- o rerank
  reordena candidatos que o 1º estágio já selecionou a partir da legenda; se a legenda
  não descreve o conteúdo real da imagem, o rerank não tem o que resgatar.

  \pause
  \begin{decisaobloco}{Decisão}
    O esforço futuro prioriza achar/aprovar um VLM de qualidade suficiente (nenhum
    testado até agora atinge o critério de 90\% de fidelidade), não otimizar o rerank
    multimodal --- que segue valioso como ganho incremental \emph{sobre} um bom
    captioning, mas não como substituto para resolver a qualidade da legenda.
  \end{decisaobloco}

  \pause
  \footnotesize \textit{Ressalva: comparação com apenas 1 VLM ruim e 1 VLM bom (ainda
  reprovado formalmente) --- evidência forte neste corpus, não prova formal universal.}
\end{frame}

% ─────────────────────────────────────────────────────────────────────
% ATO 4 — o que a régua ensinou sobre escolhas
% ─────────────────────────────────────────────────────────────────────
\begin{frame}{Ato 4 --- Escolhas que a medição sustentou}
  \small
  \begin{numerobloco}{TF-IDF vs. BM25 (agosto)}
    Empate no pipeline completo (98\%/44\%/0,866 vs. 98\%/43\%/0,866). BM25 perde no
    lexical isolado (94\% vs. 98\%). \textbf{Mantido TF-IDF.}
  \end{numerobloco}

  \begin{numerobloco}{Python puro vs. LangChain no RAG (setembro)}
    Recall@5 e MRR idênticos; Precision@5 sobe 11,4pp --- mas o ganho vem do
    \emph{algoritmo} BM25+RRF, não do framework. \textbf{Migrado assim mesmo}, por
    consistência com LangGraph nos módulos seguintes --- razão declarada, não disfarçada
    de ganho técnico.
  \end{numerobloco}

  \begin{numerobloco}{qwen2.5:7b vs. DeepSeek-R1:7b}
    2/2 sucessos em 112,6s vs. 1/2 em 258,2s (2,3x mais lento, 3x mais tokens de saída).
    \textbf{Mantido qwen2.5:7b} --- raciocínio longo não compensa em self-repair de SQL.
  \end{numerobloco}
\end{frame}

% ─────────────────────────────────────────────────────────────────────
% ATO 5 — até onde a régua alcança
% ─────────────────────────────────────────────────────────────────────
\begin{frame}{Ato 5 --- Até onde a régua alcança (e onde não alcança)}
  \small
  \begin{numerobloco}{Cobertura por categoria (vocabulário CERISE)}
    \begin{center}
    \begin{tabular}{lll}
      \toprule
      \textbf{Categoria} & \textbf{Estado} & \textbf{O que temos} \\
      \midrule
      Recuperação & \textcolor{ceriseroxo}{completa} & Recall@k, Precision@k, MRR, nDCG@k \\
      LLM         & parcial & Faithfulness, Hallucination Rate \\
      Sistema     & parcial & Tokens, latência (só nos agentes) \\
      Ferramentas & \alert{lacuna} & sem SQL/Tool Selection Accuracy formal \\
      Industrial  & \alert{lacuna} & sem Diagnosis/Root Cause Accuracy \\
      \bottomrule
    \end{tabular}
    \end{center}
  \end{numerobloco}

  \pause
  \textbf{O que falta, nomeado:} Answer Relevancy e BERTScore (categoria LLM); latência no
  dashboard de produção, não só nos agentes de teste (Sistema); e as duas categorias
  inteiras --- Ferramentas e Industrial --- que hoje não têm nenhuma métrica formal.

  \pause
  \begin{decisaobloco}{Por que isso está no documento}
    Uma apresentação que só mostra o que foi medido esconde o tamanho real do que ainda
    não sabemos. As lacunas acima são a fila de trabalho, não omissões.
  \end{decisaobloco}
\end{frame}

% ─────────────────────────────────────────────────────────────────────
% ATO 6 — a régua testada num dado que ninguém no time construiu
% ─────────────────────────────────────────────────────────────────────
\begin{frame}{Ato 6 --- A régua testada num dado que ninguém no time construiu}
  \small
  \begin{numerobloco}{O problema}
    Até 13 de setembro, toda medição de RAG multimodal do Harbor usava dado sintético
    --- 26 gráficos gerados pelo próprio time. Isso valida a arquitetura contra si
    mesma, não contra o mundo real.
  \end{numerobloco}

  \pause
  \textbf{O que fizemos:} integramos o \textbf{OpenPack} (dataset público de operações
  de embalagem logística, sensores vestíveis, 102 sessões reais) como 8º dataset do
  Harbor, aplicando \textbf{RAG-HAR} (Sivaroopan et al., arXiv:2512.08984, mai/2026)
  --- estatísticas de sensor viram texto por template determinístico, sem modelo de
  visão, sem treino.

  \pause
  \begin{decisaobloco}{A pergunta estava errada, e a régua provou isso}
    Recall@5 deu 0\% --- alarmante até investigarmos. A causa não era bug: o
    RAG-HAR nunca busca ``o documento exato'', busca vizinhos da mesma classe
    (\emph{retrieval por categoria}, não por instância). Métrica corrigida
    (k-NN label purity): \textbf{12,9\%}, contra 10\% de acaso --- sinal real,
    porém fraco. A mesma disciplina do Ato 3 se repete: quando a régua diz um
    número absurdo, a régua pode estar certa e a pergunta é que estava errada.
  \end{decisaobloco}
\end{frame}

% ─────────────────────────────────────────────────────────────────────
% Fechamento
% ─────────────────────────────────────────────────────────────────────
\begin{frame}{O que fica}
  \begin{numerobloco}{Um conjunto de decisões, três respostas diferentes}
    \begin{itemize}
      \item \textbf{Sim:} migração do self-repair e do roteador para LangGraph;
        top-p=0,7 como limiar de rerank multimodal (depois de consertada a régua);
        OpenPack como 8º dataset, integrado nas três rotas do chat sem regressão
        (roteamento 55/67 $\to$ 58/69).
      \item \textbf{Não:} Agentic RAG, Contextual Retrieval, e o rerank multimodal como
        forma de compensar captioning ruim --- todos medidos, documentados, não
        promovidos.
      \item \textbf{A pergunta estava errada, depois corrigida:} calibração de top-p
        (score=0,0 no 1º estágio) $\to$ consertada, respondida, e usada para testar uma
        pergunta melhor (o experimento decisivo do Ato 3); Recall@5=0\% do OpenPack
        $\to$ era retrieval por categoria, não por instância (Ato 6).
    \end{itemize}
  \end{numerobloco}

  \pause
  \textbf{Próximos passos:} achar/aprovar um VLM de captioning com qualidade suficiente
  (prioridade elevada pelo resultado do experimento decisivo); instrumentar trace no
  dashboard de produção; \emph{checkpointer} no agente ReAct; avaliar diagnóstico em
  camadas e NL-to-SQL (os itens de maior risco do roadmap); solicitar acesso ao RGB
  completo do OpenPack (aprovação institucional pendente) para fechar a validação
  multimodal com dado real de imagem.

  \pause
  \begin{decisaobloco}{A parte que importa}
    Boa parte dos resultados medidos nesta narrativa foi negativa --- e isso é o
    argumento mais forte deste documento. Uma régua que só confirma o que a gente já
    queria não é régua.
  \end{decisaobloco}
\end{frame}

\end{document}
